\problemsheader

\problem{Қоспа}{10.0}

\subproblem{Үйкеліс арқылы дәнекерлеу}{3.5}

\textit{Үйкеліс арқылы дәнекерлеудің} қарапайым моделін қарастырайық. Радиусы $R=\qty1{\cm}$, осьтері сәйкес келетін екі бірдей жеткілікті ұзын болат цилиндрлер берілген. Цилиндрлердің бірі қандай да бір $\omega$ бұрыштық жылдамдығымен айналдырылады және ол бекітілген цилиндрдің ұшына $p=\qty{20}{\MPa}$ тұрақты қысыммен қысылады. Шекті $\omega$-ның қандай мәнінде болат стационарлы режимде (цилиндр жеткілікті ұзақ уақыт айналғаннан кейін) балқи бастайтынын анықтаңыз. Цилиндр радиусы бойынша температура тез теңеседі деп есептеңіз. Қоршаған орта температурасы $T_0=\qty{20}{\celsius}$. Тот баспайтын болат үшін: болаттың болатқа үйкеліс коэффициенті $\mu=0.4$, жылуөткізгіштік коэффициенті $\kappa=\qty{20}{\WmK}$, мәжбүрлі конвекция кезіндегі ауаға жылу беру коэффициенті $\alpha=\qty{200}{\W\per\metre\squared\per\kelvin}$, ал балқу нүктесі $T^*=\qty{1520}{\celsius}$. Болаттың параметрлерін оның температурасына тәуелсіз деп санаңыз; цилиндрлердің жылу сыйымдылығын ескермеңіз.

\textbf{Математикалық кеңес:} Қандай да бір нақты $\lambda$ заттық тұрақтысы үшін $x''(t) = \lambda^2 x(t)$ түріндегі дифференциалдық теңдеудің шешімі келесідей болады:
$$
x(t) = Ae^{\lambda t} + Be^{-\lambda t}
$$
мұндағы $A$ және $B$ --- интегралдау тұрақтылары, ал $e=2.718...$ --- Эйлер саны.

\subproblem{Айнымалы (?) ток}{3.5}

\pic7R{0pt}{0.3}{figures/1b.1}

Оң жақтағы суретте көрсетілген тізбекте сыйымдылығы $C=\qty{4.7}{\mF}$ конденсатор мен индуктивтілігі $L=\qty{5.4}{\H}$ катушка тізбектей жалғанып, $K$ кілтіне қосылған. Кілт екі нұсқада да $U_0=\qty{10}{\V}$ тұрақты кернеу көзіне өте қысқа $\tau=\qty5{\ms}$ уақытқа қосылады, содан кейін $T=\qty{250}{\ms}$ уақыт бойы ($T\gg\tau$ деп есептеңіз) қысқа тұйықталуға, яғни ең шеткі сол жақ тік күйге ауыстырылады, содан кейін қайтадан кернеу көзіне қосылады және т.с.с. $t=0$ сәтінде екі жағдайда да тізбекте ток болған жоқ.

\newcounter{taskstar}
\setcounter{taskstar}{0}
\begin{list}{\textbf{(\alph{taskstar})}\ }{
\usecounter{taskstar}
\setlength{\leftmargin}{18pt}%
\setlength{\labelwidth}{0pt}%
\setlength{\labelsep}{0pt}%
\setlength{\itemsep}{6pt}%
\setlength{\topsep}{6pt}}
\item Конденсатор зарядталмаған, ал кілт шеткі сол жақ тік күйде болды делік. $t=0$ сәтінде кілт $U_0$-ге қосылып, жоғарыда сипатталғандай ауысып тұрды. $t=3T+\tau$ уақыт мезетіндегі ($t=3T$, $K$ кілті шеткі сол жақ күйге ауысқаннан кейінгі мезет) жүйенің энергиясы қандай болады? $0\leq t \leq 3T+\tau$ уақыт аралығындағы конденсатордың максималды заряды $q_0$ неге тең болды?
\item Айталық, $t=0$ мезетінде конденсаторда белгісіз $q_*>0$ заряды болды (полярлығы суретте көрсетілген), ал кілт шеткі сол жақ тік күйде қозғалмай тұрды. Содан кейін, қандай да бір $t=t_*>0$ мезетінде кілт $U_0$ көзіне қосылып, жоғарыда сипатталғандай ауыстырыла бастады. $K$ кілтін ауыстырып қосқан кезде жүйенің толық энергиясы өзгеріссіз қалғаны белгілі болды. Бұл $q_*$ және $t_*$-ның қандай мәндерінде мүмкін?
\end{list}

\subproblem{Дыбыс интерференциясы}{3.0}

Екі шексіз қатты жазықтық (қабырғалар) тік бұрыш жасап қиылысады. Шексіздіктен жазық дыбыс толқыны келіп түседі, оның толқындық векторы дәл жазықтықтар арасындағы бұрыштың биссектрисасы бойымен бағытталған. Түскен толқындағы ауа молекулалары жылдамдығының амплитудасы $v$, ал толқын ұзындығы $\lambda$-ға тең. Қабырғалар идеал шағылыстырады (шағылысу энергия шығынынсыз жүреді); $\lambda$ ретіндегі қашықтықтағы дифракциялық эффектілерді ескермеңіз.

\setcounter{taskstar}{0}
\begin{list}{\textbf{(\alph{taskstar})}\ }{
\usecounter{taskstar}
\setlength{\leftmargin}{18pt}%
\setlength{\labelwidth}{0pt}%
\setlength{\labelsep}{0pt}%
\setlength{\itemsep}{6pt}%
\setlength{\topsep}{6pt}}
\item Осы жүйедегі ауа қозғалысының мүмкін болатын максималды жылдамдығын табыңыз.
\item Кеңістіктің қай нүктелерінде осы максималды жылдамдыққа қол жеткізіледі?
\end{list}